\section{Tema y objetivo del proyecto}
\label{sec:tema}

\subsection{Título}
\textbf{Simulador de Entrevistas Laborales Adaptativo}: un reclutador basado en
LLM que adapta sus preguntas a la industria y nivel del candidato, analizando
en tiempo real la voz (fluidez, muletillas, pausas, tono) y el lenguaje
corporal por webcam (contacto visual, postura, gestualidad), y que habilita
\textit{mock interviews} colaborativas entre estudiantes con roles
intercambiables.

\subsection{Objetivo principal}
Diseñar y construir una aplicación web/PWA que permita a estudiantes y
recién egresados practicar entrevistas laborales en cualquier momento,
recibiendo retroalimentación objetiva, instantánea y multimodal sobre su
desempeño verbal y no verbal, así como un plan de mejora personalizado y la
opción de validación entre pares.

\subsection{Objetivos específicos}
\begin{enumerate}
  \item Implementar un agente entrevistador conversacional por voz capaz de
        adaptar dinámicamente la dificultad y el dominio (industria, rol,
        nivel) según el perfil del candidato.
  \item Calcular y mostrar en tiempo real métricas de comunicación verbal
        (palabras-muletilla por minuto, ritmo, claridad) y no verbal
        (contacto visual, postura, gestualidad) sin enviar video crudo a
        servicios externos.
  \item Generar un plan de mejora individualizado tras cada sesión, basado
        en debilidades recurrentes detectadas a lo largo del historial
        longitudinal del usuario.
  \item Habilitar sesiones \textit{peer-to-peer} entre estudiantes mediante
        WebRTC, con panel de observador, comentarios anclados a
        \textit{timestamps} e intercambio de roles.
  \item Aplicar elementos de gamificación (rangos por competencia, ligas
        temáticas, \textit{badges}) que sostengan la motivación intrínseca
        sin recurrir a \textit{dark patterns}.
  \item Garantizar accesibilidad WCAG 2.2 nivel AA y rutas alternativas para
        usuarios con discapacidad visual, auditiva, motora o cognitiva.
\end{enumerate}

\subsection{Problema real que resuelve}
Los estudiantes y recién egresados enfrentan dos barreras al prepararse para
entrevistas técnicas:

\begin{itemize}
  \item \textbf{Falta de retroalimentación objetiva}: amigos y familiares no
        pueden evaluar lenguaje corporal y fluidez de manera consistente, y
        el sesgo afectivo distorsiona su juicio.
  \item \textbf{Costo y escasez de coaches profesionales}: el coaching uno
        a uno es caro, requiere agenda y suele ser inaccesible para
        estudiantes universitarios.
\end{itemize}

\noindent Una herramienta multimodal con IA permite practicar
$24\times 7$, sin exposición inicial a juicio humano, y luego con
validación opcional entre pares.

\subsection{Diferenciación frente a propuestas similares}
A diferencia de aplicaciones para entrenamiento de habilidades sociales
orientadas al bienestar emocional, esta propuesta:

\begin{itemize}
  \item Se enfoca en empleabilidad técnica (entrevistas reales de trabajo).
  \item Usa análisis audiovisual multimodal en tiempo real (voz $+$ video),
        no solo diálogo textual.
  \item Incorpora una dimensión comunitaria y competitiva (peer mock,
        ligas, \textit{ranking}).
  \item Su métrica objetivo es desempeño profesional, no reducción de
        síntomas clínicos.
\end{itemize}
